\section{Experiments}
\label{sec:exp}

\subsection{Data and setup}
\label{sec:data}

We use two frozen public snapshots; no model is trained or evaluated by us. \textbf{HELM MMLU} (release v1.13.0, 2025-01-10): the per-subject exact-match tables of 57 MMLU subjects, with all $M=79$ models present on every subject ($n=57$). \textbf{Open LLM Leaderboard} (\texttt{contents}, commit \texttt{9c09a7ca}, last updated 2025-03-20): $M=4497$ models after dropping 79 duplicated names (keeping the best-scoring variant) and $n=6$ task families (IFEval, BBH, MATH Lvl~5, GPQA, MuSR, MMLU-Pro) on the leaderboard's normalised $0$--$100$ scale; one flagged model is removed, and 708 of the models are merges (kept in the main analysis, dropped in a robustness slice). The two datasets span two regimes: few tasks and many models, and many correlated tasks and few models. Neither source declares a data licence (the HELM code is Apache-2.0); we use public aggregate scores only, redistribute only derived tables and code, and give versions, checksums and our licence judgement in the reproduction notes. The analysis runs on CPU only and trains no model.

Every model whose rank under the mean is at least 2 is a target (HELM: 78 targets; Open LLM Leaderboard: 4495 for the score rules). Ordinal rules on the Open LLM Leaderboard are run on a rank-stratified sample of targets (fixed seed, drawn once) because their mixed-integer programs are costly; for ranks 2--100 they were also run on every target (Appendix~\ref{app:sens}). Each target and (rule, cap) pair is solved exactly with the solvers of Section~\ref{sec:method}, margin $\varepsilon=10^{-3}$. Targets are the whole population of each snapshot, so intervals over targets are descriptive: they are 95\% percentile bootstrap intervals over the targets in a bin, and Wilson intervals only for the shares in Figure~\ref{fig:cap}. They cover the sampling of targets within one snapshot and nothing else: not the choice of snapshot, of task set or of submitted models, not evaluation noise (Section~\ref{sec:exp-robust}), and not the dependence between rows, since many models are fine-tunes or merges of one another. They should be read as a measure of spread, not of the precision of a population parameter.

\paragraph{Empirical gains.}
The empirical caps use the 540 (fine-tune, base) pairs of the Open LLM Leaderboard for which the model's \texttt{Base Model} field names another unflagged model of the leaderboard (the field is user-entered metadata and we did not audit it); the gain on a task is the fine-tune's score minus the base's. The pairs mix instruction and chat tunes, merges and domain fine-tunes and are not outcomes of targeted training; they are a proxy with unknown selection. A restricted set without chat-typed models, merges and instruction-named tunes (\pairRestrictedN{} pairs) gives a pooled 95\% quantile of \empPooledRestricted{} instead of \empPooled{} points and nearly the same infeasible shares (Appendix~\ref{app:sens}).

\subsection{Closed forms and rule structure}
\label{sec:exp-theory}

\paragraph{Solvers match the closed forms.}
On the 237 HELM rows (mean, minimum, median; ceiling cap) the solvers and Propositions~\ref{prop:mean}--\ref{prop:median} agree to $|\Delta L_1|\le\thMaxLone$, $|\Delta K|\le\thMaxK$ and $|\Delta L_\infty|\le\thMaxLinf$ (bisection tolerance; Table~\ref{tab:theory}, appendix). Under the ceiling the equality $L_1^\ast=nL_\infty^\ast$ of Proposition~\ref{prop:mean} applies to only 4 of 79 HELM targets, because a uniform raise clips at 100 for most; in general $L_1^\ast\le nL_\infty^\ast$, with equality when no cap binds.

\paragraph{Minimum versus mean.}
For every target with $L_1>0$ (\ctwoHelmShare\% of the \ctwoHelmN{} HELM targets and \ctwoOllbShare\% of the \ctwoOllbN{} Open LLM Leaderboard targets) the minimum rule needs fewer points than the mean, so the expectation that the minimum is dearer for balanced targets is contradicted. The median ratio $L_1^{\min}/L_1^{\mathrm{mean}}$ is \ctwoHelmRatio{} on HELM (95\% interval \ctwoHelmRatioLo--\ctwoHelmRatioHi) and \ctwoOllbRatio{} on the Open LLM Leaderboard (\ctwoOllbRatioLo--\ctwoOllbRatioHi). With the dispersion of a target's scores measured by the coefficient of variation (CV) of its task scores, the ratio \emph{rises} with CV on both datasets: the median ratio in the low, middle and high CV terciles is \ctwoHelmTLow, \ctwoHelmTMid{} and \ctwoHelmTHigh{} on HELM (Spearman $\rho=$~\ctwoHelmRho) and \ctwoOllbTLow, \ctwoOllbTMid{} and \ctwoOllbTHigh{} on the Open LLM Leaderboard ($\rho=$~\ctwoOllbRho): the minimum saves most for balanced targets. This is specific to the statistic: with the spread measured as mean minus minimum, $\rho$ is \ctwoHelmRhoMM{} on HELM and \ctwoOllbRhoMM{} on the Open LLM Leaderboard, so the dependence on dispersion is not stable across statistics and datasets. Cheaper in points does not mean more often first: under additive gain caps the minimum is the least reachable rule on HELM (Section~\ref{sec:exp-caps}), so cost and feasibility must be read together.

\begin{figure}[tb]
\centering
\includegraphics[width=\linewidth]{../figures/fig_key_results}
\caption{(a) Cost of the minimum rule relative to the mean, $L_1^{\min}/L_1^{\mathrm{mean}}$ at cap 100: median per tercile of the coefficient of variation (CV) of the target's scores; horizontal bar the tercile's CV range, vertical bar the 95\% bootstrap interval over targets; HELM filled (78 targets), Open LLM Leaderboard open (4495); Spearman $\rho$ over all targets is printed. (b) Open LLM Leaderboard, targets ranked 2--100: median $L_1$ of the median rule and of the minimum rule on four data slices (main sample, no merged models, raw scores, top 500) with 95\% bootstrap intervals and $n$ targets; compare the order within a slice, not levels across slices.}
\label{fig:key}
\end{figure}

\paragraph{Task-subset scaling.}
We draw random task subsets of $n=5,\dots,57$ from HELM (50 draws, cap $+10$, two targets of rank 2 and 5, so this is a small study; Figure~\ref{fig:tasksub}). In $K/n$ at $n=57$ the mean needs \cthreeFullMean, Borda \cthreeFullBorda, plurality \cthreeFullPlur, the median \cthreeFullMedian{} and the minimum \cthreeFullMin. The median is above the minimum at every $n$ tested, so the expectation that the median is cheaper than the minimum in tasks touched was not borne out; at $n=5$ the minimum is infeasible for \cthreeMinFeasFive{} of the draws of the rank-5 target.

\subsection{Cost versus current rank}
\label{sec:exp-rank}

\begin{figure}[tb]
\centering
\includegraphics[width=0.72\linewidth]{../figures/fig_rank_rules}
\caption{Minimum investment to reach first place versus current rank under the mean (cap 100). Top: tasks touched $K$ (mean per rank bin); bottom: points $L_1$ (median per bin, symlog axis). Bands: 95\% bootstrap intervals over the targets in a bin. Left: Open LLM Leaderboard; right: HELM MMLU. $K$-minimal and $L_1$-minimal solutions are separate optima.}
\label{fig:rank}
\end{figure}

Figure~\ref{fig:rank} shows that on both datasets cost grows with rank, with rule-specific slopes. For the mean the dependence is an identity: $L_1^\ast$ is linear in the gap to the best rival (Proposition~\ref{prop:mean}), so $R^2=1$ against the gap is not evidence. For the geometric and harmonic means and the median the cost is also close to a function of the gap in rule units by construction (the rule value must rise by the gap), so the informative quantity is what the gap adds to rank. We regress, for every rule but the mean, $\log L_1$ on $\log$(gap in rule units) and on $\log$(rank) (ceiling cap, targets with positive gap and cost) and report the difference of the two $R^2$ with a bootstrap interval over targets (Table~\ref{tab:c4}, appendix). The gap beats rank, with an interval for the $R^2$ difference that excludes zero, for the geometric mean ($R^2$ \cfourHelmGeomeanGap{} vs.\ \cfourHelmGeomeanRank{} on HELM, \cfourOllbGeomeanGap{} vs.\ \cfourOllbGeomeanRank{} on the Open LLM Leaderboard), the harmonic mean and the median on both datasets, and for Borda on HELM (\cfourHelmBordaGap{} vs.\ \cfourHelmBordaRank{}); on the Open LLM Leaderboard the margins are small (0.02--0.03). It is \emph{not} supported for the minimum (difference \cfourHelmMinDiff{} [\cfourHelmMinDiffLo, \cfourHelmMinDiffHi] on HELM, \cfourOllbMinDiff{} [\cfourOllbMinDiffLo, \cfourOllbMinDiffHi] on the Open LLM Leaderboard), the Condorcet requirement (\cfourHelmCondorcetDiff{} [\cfourHelmCondorcetDiffLo, \cfourHelmCondorcetDiffHi]; \cfourOllbCondorcetDiff{} [\cfourOllbCondorcetDiffLo, \cfourOllbCondorcetDiffHi]), Borda on the Open LLM Leaderboard or plurality on HELM; Table~\ref{tab:c4} gives all values. Plurality on the Open LLM Leaderboard cannot be tested because rank ties leave too few distinct costs, and the ordinal rules there use a rank-stratified sample, so their intervals are wide. Rank alone is already a strong predictor everywhere ($R^2\ge0.72$, and $\ge0.84$ for all score rules), so the gap adds explanatory power only where stated above.

\subsection{Caps decide feasibility}
\label{sec:exp-caps}

\paragraph{Frontier caps make first always reachable.}
Under the ceiling, front+0 and front+2, first place is feasible for every target under every rule on both datasets. This is mechanical: a target raised to the best rival score on every task weakly dominates every rival everywhere, so it is first under any monotone rule. The frontier caps therefore cannot separate rules in feasibility; they only change cost.

\paragraph{Fixed per-task gains separate rules, differently on each dataset.}
With cap $+10$ the share of HELM targets that can reach first ranges from \gfTenHelmPlurality\% (plurality) through \gfTenHelmMedian\% (median), \gfTenHelmMean\% (mean) and \gfTenHelmBorda\% (Borda) to \gfTenHelmGeomean\% (geometric), \gfTenHelmHarmonic\% (harmonic) and \gfTenHelmMin\% (minimum; Condorcet \gfTenHelmCondorcet\%). On the Open LLM Leaderboard (all 4495 targets, score rules) it is \gfTenOllbMean\% (mean), \gfTenOllbMedian\% (median), \gfTenOllbGeomean\% (geometric), \gfTenOllbHarmonic\% (harmonic) and \gfTenOllbMin\% (minimum): the minimum is the \emph{least} reachable rule on HELM and the \emph{most} reachable of these five on the Open LLM Leaderboard. The ordinal rules there were run on all \cvxOllbLinearN{} targets of ranks 2--100 only, where the shares at cap $+10$ are plurality \ordHundredTenPluralityFeas\%, Condorcet \ordHundredTenCondorcetFeas\%, Borda \ordHundredTenBordaFeas\% against \ordHundredTenMeanFeas\% for the mean, \ordHundredTenMedianFeas\% for the median and \ordHundredTenMinFeas\% for the minimum; these shares are for the top of the board and must not be compared with the all-target figures. Their minimum uniform gain cap $h^\ast$ has median \hstarOrdBordaElevenThirty{} (Borda), \hstarOrdCondorcetElevenThirty{} (Condorcet) and \hstarOrdPluralityElevenThirty{} (plurality) points for ranks 11--30, against \hstarBinOllbMeanElevenThirty{} for the mean, and at most \hstarOrdPluralityMax{} points in ranks 2--100, below the pooled empirical quantile (\empPooled).

\begin{figure}[tb]
\centering
\includegraphics[width=0.8\linewidth]{../figures/fig_cap_curve}
\caption{Share of targets with $h^\ast\le h$ (95\% Wilson intervals). (a) HELM, 78 targets, all rules. (b) Open LLM Leaderboard, all 4495 targets, score rules only. (c) Open LLM Leaderboard, all 99 targets of ranks 2--100, all six rules on the same targets (ordinal rules were solved for these only). Dotted: pooled empirical 95th percentile of fine-tuning gains (25.3 points); grey: its range over the six tasks (3.2--46.0); dashed (HELM): the MMLU-Pro value (3.2).}
\label{fig:cap}
\end{figure}

\paragraph{Minimum cap $h^\ast$ and the empirical caps.}
The minimum uniform gain cap $h^\ast$ removes the choice of a few scenarios. It depends strongly on the target's rank, so Table~\ref{tab:hstar} reports it by rank bin. Under the mean, the median $h^\ast$ is \hstarBinHelmMeanElevenThirty{} points on HELM and \hstarBinOllbMeanElevenThirty{} points on the Open LLM Leaderboard for ranks 11--30, and \hstarBinHelmMeanThirtyOneHundred{} and \hstarBinOllbMeanThirtyOneHundred{} for ranks 31--100; the \hstarBinNOllbOverHundred{} Open LLM Leaderboard targets ranked beyond 100 (\ollbShareOverHundred\% of all) need \hstarBinOllbMeanOverHundred{}. The all-target medians, \hstarHelmMean{} points on HELM (95\% interval \hstarHelmMeanLo--\hstarHelmMeanHi) and \hstarOllbMean{} on the Open LLM Leaderboard (\hstarOllbMeanLo--\hstarOllbMeanHi), are dominated by low-ranked targets and should not be read as the investment of a contender for first place. Across rules the all-target medians run on HELM from \hstarHelmPlurality{} (plurality) to \hstarHelmMin{} (minimum), with the mean at \hstarHelmMean, and on the Open LLM Leaderboard from \hstarOllbMin{} (minimum) to \hstarOllbMedian{} (median), with the mean at \hstarOllbMean: the order of the minimum and the median reverses between datasets, and on the Open LLM Leaderboard the order of their $L_1$ also depends on the data slice (Appendix~\ref{app:robust}). We compare with observed fine-tuning gains: the pooled 95th percentile of 540 fine-tune-minus-base differences on the Open LLM Leaderboard is \empPooled{} points, and the per-task 95th percentiles range from 3.2 to 46.0. Under the pooled value the share of targets whose $h^\ast$ exceeds it is \hstarShareHelmMean\% (mean, HELM) and \hstarShareOllbMean\% (mean, Open LLM Leaderboard); for the minimum it is \hstarShareHelmMin\% and \hstarShareOllbMin\%. Under the \emph{per-task} 95th-percentile caps, which are far more restrictive, first place is unreachable for \ceightInfHelmMean\% of HELM targets under the mean (\ceightInfHelmPlurality\% under plurality) and for \ceightInfOllbMean\% of Open LLM Leaderboard targets under the mean (\ceightInfOllbMin\% under the minimum). On HELM these caps are transferred from the MMLU-Pro quantile (\empMmluQ{} points, against \empPooled{} for the pooled scalar) because HELM has no base/fine-tune pairs; with the pooled scalar instead the feasible share is \empPoolHelmMin--\empPoolHelmMax\%, so the HELM conclusion depends strongly on this transfer. Under the MMLU-Pro value the share of HELM targets whose $h^\ast$ exceeds the cap is \hstarMmluShareHelmPlurality--\hstarMmluShareHelmHarmonic\% across rules, against \hstarShareHelmMean\% (mean) under the pooled scalar.

\begin{table}[tb]
\caption{Minimum uniform gain cap $h^\ast$ (points) for first place by rank bin: median with 95\% bootstrap interval over targets (none for fewer than 3 targets) for the mean, median and minimum rules, and the share of targets whose $h^\ast$ exceeds the pooled 95th percentile of fine-tuning gains (\empPooled) and, on HELM, the MMLU-Pro value that the HELM empirical caps use (\empMmluQ). Infeasible counts as exceeding. \ollbShareOverHundred\% of the Open LLM Leaderboard targets rank beyond 100 and dominate its all-target median (\hstarOllbMean); rank 2 is one target. Score rules only.}
\label{tab:hstar}
\centering\scriptsize
\resizebox{\linewidth}{!}{\input{../figures/tables/tab_hstar_bins}}
\end{table}

\subsection{Ordinal rules}
\label{sec:exp-ordinal}

Table~\ref{tab:main} (appendix) gives the ceiling, cap $+10$ and empirical-95\% blocks. On HELM, rank bin 11--30, ceiling cap, median $K$ and $L_1$ are: plurality \binKHelmPlurality{} tasks and \binLHelmPlurality{} points; mean \binKHelmMean{} and \binLHelmMean; Condorcet \binKHelmCondorcet{} and \binLHelmCondorcet; Borda \binKHelmBorda{} and \binLHelmBorda. Borda and the Condorcet requirement therefore need many tasks, but, contrary to our expectation, the Condorcet requirement is \emph{cheaper} than the mean in points; plurality is cheapest in both. Because the Condorcet requirement is sufficient for first under any Condorcet-consistent rule, these numbers are an upper bound for Copeland, Schulze or Kemeny. On the Open LLM Leaderboard, over all \cvxOllbLinearN{} targets of ranks 2--100 (ceiling), Borda, plurality and the Condorcet requirement touch a median of \ordHundredCeilBordaK, \ordHundredCeilPluralityK{} and \ordHundredCeilCondorcetK{} tasks, but under cap $+10$ plurality is the hardest rule to make feasible (\ordHundredTenPluralityFeas\% of targets, against \ordHundredTenBordaFeas\% for Borda and \ordHundredTenCondorcetFeas\% for the Condorcet requirement; Appendix~\ref{app:sens}).

\subsection{Robustness}
\label{sec:exp-robust}

Each draw needs an exact solve, so the noise and task-bootstrap runs use rank-stratified samples of \noiseNHelm{} HELM and \noiseNOllb{} Open LLM Leaderboard targets (\noiseDrawsHelm{} and \noiseDrawsOllb{} noise draws, \tbootDrawsHelm{} bootstrap draws; score rules only on the latter). Under evaluation noise (binomial standard errors from item counts, cap $+10$) the ratio of $L_1$ to its observed value (median over targets, mean over draws) is \noiseRatioHelmMean{} for the mean and \noiseRatioHelmMin{} for the minimum on HELM, \noiseRatioOllbMean{} and \noiseRatioOllbMin{} on the Open LLM Leaderboard; the minimum rule can reach first place in only \noiseFeasHelmMin\% of noisy HELM target-draws, against \noiseFeasObsHelmMin\% of the targets on the observed scores. A target that is first by exactly $\varepsilon=10^{-3}$ stays first with probability only about \flipPHelmMean{} (mean), \flipPHelmMedian{} (median), \flipPHelmMin{} (minimum) on HELM (ranks 11--30), so $\varepsilon$ is not a robust margin. Margins of $10^{-6}$ and $0.1$ change the median $L_1$ by at most a few per cent, but rule-specific margins halve the minimum's advantage over the mean and collapse its feasibility under fixed gains (Appendix~\ref{app:sens}); data slices change the order of the median and the minimum rule's cost (Appendix~\ref{app:robust}).
